\begin{table}[t]
\centering
\small
\setlength{\tabcolsep}{6pt}
\begin{tabular}{llcc}
\toprule
System & Comparison & Accuracy & AUC \\
\midrule
\multicolumn{4}{l}{\textit{S vs C1 (same-pipeline cover)}} \\
TopicQA   & S vs C1 & $0.490 \pm 0.024$ & $0.355 \pm 0.059$ \\
StorySlot & S vs C1 & $0.480 \pm 0.015$ & $0.401 \pm 0.052$ \\
LitReview & S vs C1 & $0.493 \pm 0.010$ & $0.365 \pm 0.037$ \\
\midrule
\multicolumn{4}{l}{\textit{S vs C2 (prompted cover, length-matched)}} \\
TopicQA   & S vs C2 & $0.882 \pm 0.101$ & $0.945 \pm 0.060$ \\
StorySlot & S vs C2 & $0.972 \pm 0.012$ & $0.996 \pm 0.002$ \\
LitReview & S vs C2 & $0.825 \pm 0.041$ & $0.902 \pm 0.037$ \\
\midrule
\multicolumn{4}{l}{\textit{S vs C3 (harder cover, StorySlot only)}} \\
StorySlot & S vs C3 & $0.698 \pm 0.050$ & $0.748 \pm 0.056$ \\
\bottomrule
\end{tabular}
\caption{Fine-tuned DistilBERT-base binary classifier, 5-fold stratified CV.
S vs C1 is a clean null across systems (accuracy at 50\%, AUC below chance
indicating folds with inverted predictions on a non-existent signal). S vs
C2 is highly detectable for all three systems --- the classifier picks up
on length-matched prompted-cover artifacts (formatting, register), not on
the steganographic channel itself. S vs C3 (a harder, distribution-matched
cover available only for StorySlot) reduces detectability substantially
($0.97 \to 0.70$), supporting the C2 artifact interpretation. C3 was not
constructed for TopicQA or LitReview in this batch.}
\label{tab:steg-transformer}
\end{table}